\documentclass[
	% -- opções da classe memoir --
	article, % indica que é um artigo acadêmico
	11pt, % tamanho da fonte
	oneside, % para impressão apenas no recto. Oposto a twoside
	a4paper, % tamanho do papel.
	% -- opções da classe abntex2 --
	%chapter=TITLE,		% títulos de capítulos convertidos em letras maiúsculas
	%section=TITLE,		% títulos de seções convertidos em letras maiúsculas
	%subsection=TITLE,	% títulos de subseções convertidos em letras maiúsculas
	%subsubsection=TITLE % títulos de subsubseções convertidos em letras maiúsculas
	% -- opções do pacote babel --
	english, % idioma adicional para hifenização
	brazil, % o último idioma é o principal do documento
	sumario=tradicional
]{abntex2}

% ---
% PACOTES
% ---

% ---
% Pacotes fundamentais
% ---
\usepackage{lmodern}			% Usa a fonte Latin Modern
\usepackage[T1]{fontenc}		% Selecao de codigos de fonte.
\usepackage[utf8]{inputenc}		% Codificacao do documento (conversão automática dos acentos)
\usepackage{indentfirst}		% Indenta o primeiro parágrafo de cada seção.
\usepackage{nomencl} 			% Lista de simbolos
\usepackage{color}				% Controle das cores
\usepackage{graphicx}			% Inclusão de gráficos
\usepackage{microtype} 			% Para melhorias de justificação
\usepackage{hyperref}           % Habilita links e referências

% --- Header e Footer ---

\setlength{\headheight}{36pt}
\setlength{\headsep}{10pt}

\makepagestyle{meuestilo}
\makeevenhead{meuestilo}{\includegraphics[height=1.1cm]{FCI.png}}{}{\thepage}
\makeoddhead {meuestilo}{\includegraphics[height=1.1cm]{FCI.png}}{}{\thepage}
\makeheadrule{meuestilo}{\textwidth}{0.4pt}

\pagestyle{meuestilo}
% Para páginas "plain" (como início de capítulo/título):
\aliaspagestyle{chapter}{meuestilo}
\aliaspagestyle{plain}{meuestilo}

% Pacotes de citações
% ---
%\usepackage[brazilian,hyperpageref]{backref}	 % Paginas com as citações na bibl
\usepackage[alf]{abntex2cite} % Citações padrão ABNT
% ---

% ---
% Configurações do pacote backref
% Usado sem a opção hyperpageref de backref
%\renewcommand{\backrefpagesname}{Citado na(s) página(s):~}
% Texto padrão antes do número das páginas
%\renewcommand{\backref}{}
% Define os textos da citação
%\renewcommand*{\backrefalt}[4]{
%	\ifcase #1 %
%		Nenhuma citação no texto.%
%	\or
%		Citado na página #2.%
%	\else
%		Citado #1 vezes nas páginas #2.%
%	\fi}%
% ---

% --- Informações de dados para CAPA e FOLHA DE ROSTO ---

\titulo{Implementação de Ontologia para Gêmeos Digitais Urbanos}
\def \theforeigntitle{}

\autor{Daniel Yoshio Omiya \textsuperscript{1}
\and Gabriel Mason Guerino \textsuperscript{1}
\and Pedro Bacic Branco \textsuperscript{1}
\and Prof. Dr. Fabio Silva Lopes \textsuperscript{1,2} \\
\\
\textsuperscript{1}Faculdade de Computação e Informática (FCI) \\ Universidade Presbiteriana Mackenzie São Paulo, SP – Brasil \\
\\ \textsuperscript{2}Programa de Pós-Graduação em Computação Aplicada \\ Universidade Presbiteriana Mackenzie São Paulo, SP – Brasil \\
\\
\url{{10356170,10409928,10382476}@mackenzista.com.br}
\\
\url{fabio.lopes@mackenzie.br}}

\data{2026}
% ---

% ---
% Configurações de aparência do PDF final

% alterando o aspecto da cor azul
\definecolor{blue}{RGB}{41,5,195}

% informações do PDF
\makeatletter
\let\@fnsymbol\@arabic
\hypersetup{
     	%pagebackref=true,
		%pdftitle={\@title}, 
		%pdfauthor={\@author},
    	%pdfsubject={Modelo de artigo científico com abnTeX2},
	    %pdfcreator={LaTeX with abnTeX2},
		%pdfkeywords={abnt}{latex}{abntex}{abntex2}{atigo científico}, 
		colorlinks=true,       		% false: boxed links; true: colored links
    	linkcolor=blue,          	% color of internal links
    	citecolor=blue,        		% color of links to bibliography
    	%filecolor=magenta,         % color of file links
		urlcolor=blue,
		%bookmarksdepth=4
}
\makeatother
% ---

% ---
% compila o indice
% ---
\makeindex
% ---

% ---
% Altera as margens padrões
% ---
\setlrmarginsandblock{3cm}{3cm}{*}
\setulmarginsandblock{3cm}{3cm}{*}
\checkandfixthelayout
% ---

% ---
% Espaçamentos entre linhas e parágrafos
% ---

% O tamanho do parágrafo é dado por:
\setlength{\parindent}{1.3cm}

% Controle do espaçamento entre um parágrafo e outro:
\setlength{\parskip}{0.2cm} % tente também \onelineskip

% Espaçamento simples
\SingleSpacing

\begin{document}
\selectlanguage{brazil}

% Retira espaço extra obsoleto entre as frases.
\frenchspacing

% ----------------------------------------------------------
% ELEMENTOS PRÉ-TEXTUAIS
% ----------------------------------------------------------

%---
%
% Se desejar escrever o artigo em duas colunas, descomente a linha abaixo
% e a linha com o texto ``FIM DE ARTIGO EM DUAS COLUNAS''.
% \twocolumn[    		% INICIO DE ARTIGO EM DUAS COLUNAS
%
%---

% página de titulo principal (obrigatório)
\maketitle
% titulo em outro idioma (opcional)

\begin{resumoumacoluna}
  \textit{Plataformas de Gêmeos Digitais Urbanos (GDU) integram dados provenientes de múltiplos sistemas heterogêneos para suportar o monitoramento e a gestão inteligente de ambientes urbanos.
    No entanto, a ausência de um vocabulário semântico comum entre esses sistemas compromete a
    interoperabilidade e limita a capacidade de raciocínio automatizado sobre os dados coletados.
    Este trabalho propõe a modelagem e implementação de uma ontologia de domínio, alinhada aos padrões
    CityGML e SOSA/SSN, para representar entidades físicas, sensores e relações espaciais no contexto
    de um campus universitário (tomado como instância controlada de GDU).
    A pesquisa adota a metodologia Design Science Research (DSR), compreendendo a revisão dos padrões
    de referência, a modelagem da ontologia, a definição de um mecanismo de exposição, e a validação.
    Os resultados esperados incluem uma ontologia funcional capaz de responder a perguntas semânticas
    do campus, promovendo interoperabilidade semântica entre fontes de dados diversas na plataforma GDU
    do campus Higienópolis da Universidade Presbiteriana Mackenzie.
  }

  \vspace{\onelineskip}

  \noindent
  \textbf{Palavras-chave}: cidades inteligentes; gêmeos digitais urbanos; ontologia; dados semânticos; CityGML; SOSA/SSN.
\end{resumoumacoluna}

\renewcommand{\resumoname}{Abstract}
\begin{resumoumacoluna}
  \begin{otherlanguage*}{english}
    \textit{Urban Digital Twin (UDT) platforms integrate data from multiple heterogeneous systems to support the monitoring and intelligent management of urban environments.
      However, the absence of a common semantic vocabulary among these systems compromises
      interoperability and limits the ability to perform automated reasoning over the collected data.
      This work proposes the modeling and implementation of a domain ontology, aligned with the CityGML
      and SOSA/SSN standards, to represent physical entities, sensors, and spatial relationships within
      the context of a university campus (used as a controlled UDT instance).
      The research adopts the Design Science Research (DSR) methodology, encompassing the review of
      reference standards, ontology modeling, the definition of an exposure mechanism, and validation.
      The expected results include a functional ontology capable of answering semantic queries about the
      campus, promoting semantic interoperability among diverse data sources within the UDT platform of
      Universidade Presbiteriana Mackenzie.
    }

    \vspace{\onelineskip}

    \noindent
    \textbf{Keywords}: smart cities; urban digital twins; ontology; semantic data; CityGML; SOSA/SSN.
  \end{otherlanguage*}
\end{resumoumacoluna}

% ]  				% FIM DE ARTIGO EM DUAS COLUNAS
% ---

% ----------------------------------------------------------
% ELEMENTOS TEXTUAIS
% ----------------------------------------------------------
\textual

% ----------------------------------------------------------
% Introdução
% ----------------------------------------------------------
\section{Introdução}

De acordo com o relatório de Perspectiva de Urbanização Mundial da \citeonline{undesa2025wup}, as
cidades são o lar de mais de 45\% da população mundial, dobro da proporção de 1950, e projetado
para chegar a dois terços até 2050.
Esse fenômeno de crescimento acelerado das cidades eleva grandemente a complexidade de prover
infraestrutura de transporte, energia, saneamento, e governança; torna ineficiente e insuficiente
os modelos tradicionais de gestão urbana, muitas vezes com dados fragmentados em diversos sistemas
heterogêneos.
Nesse cenário, surgem as \textit{Smart Cities} (Cidades Inteligentes), dotadas do uso de Internet
das Coisas (IoT), de Inteligência Artificial (IA), e da análise de grandes quantidades de dados
(Big Data), que representaram um salto no planejamento urbano, melhorando a qualidade de vida, a
eficiência da operação e dos serviços urbanos, e a tomada de decisão pública \cite{WEIL2023104862}.

Nesse contexto, os Gêmeos Digitais Urbanos (GDU) despontam como uma tecnologia para evoluir ainda
mais a gestão e o planejamento urbano.
Sendo o Gêmeo Digital definido como uma representação virtual de um objeto da realidade, atualizado
constantemente num fluxo de dados bidirecional entre os mundos físico e digital.
Quando aplicado à escala urbana, torna-se uma plataforma capaz de monitorar, analisar e prever
sistemas complexos em tempo real.
Casos como o projeto \textit{Virtual Singapore} e a iniciativa Helsinki 3D+ demonstram o poder
deste tipo de integração, sendo usados até mesmo para resposta a emergências
\cite{FERREBIGORRA2022103905, WEIL2023104862}.

Embora seja uma tecnologia com um grande potencial, ela enfrenta uma série de desafios na sua
concepção, oriundos ainda da fragmentação dos dados necessários em sistemas que falam um dialeto
diferente cada, como: sensores de IoT, modelos de informação de edificações, dados geoespaciais em
padrões como CityGML, imagens de câmeras, drones e estações meteorológicas, entre outros.
Por exemplo, um sensor de IoT pode registrar a temperatura como um valor inteiro em graus Célsius,
enquanto a estação meteorológica poderia usar Fahrenheit.
Essa dificuldade de interoperabilidade é chamada de \textit{semantic gap} (lacuna semântica), como
apontado por \citeonline{WEIL2023104862}.

Superar essa dificuldade requer revisar o conhecimento em comum, atribuindo um significado formal e
objetivo a cada entidade presente no domínio.
Nesse sentido, as ontologias formais apresentam-se como uma abordagem promissora
\cite{metral_ontologybased_2010}.
Elas permitem que sistemas compreendam e interoperem com dados de múltiplas origens a partir de um
vocabulário comum e de regras de inferência lógica.

O presente trabalho tem como objeto de estudo a aplicação dessas tecnologias no contexto do campus
Higienópolis da Universidade Presbiteriana Mackenzie (UPM), adotando como instância representativa
um GDU em escala reduzida.
A pesquisa é conduzida sob o paradigma da \textit{Design Science Research} (DSR), que orienta o
processo de construção e avaliação de artefatos computacionais como resposta a problemas práticos
relevantes \cite{hevner_2004}, sendo a ontologia de domínio proposta o artefato central a ser
desenvolvido, avaliado e documentado.

O restante deste trabalho está organizado da seguinte forma: o Capítulo 2 apresenta a revisão
teórica dos fundamentos e padrões envolvidos; o Capítulo 3 descreve a metodologia adotada; o
Capítulo 4 detalha a arquitetura e a modelagem da ontologia proposta; o Capítulo 5 apresenta a
validação por casos de uso; e o Capítulo 6 traz as considerações finais.

\section{Referencial Teórico}

O Gêmeo Digital é um conceito introduzido por Michael Grieves em 2003, quando "apresentou a versão
digital de um produto físico" \apud[p.~7]{Opoku_2021}{Gomes_2025}, e trata-se desta representação
virtual de um objeto real que é atualizada continuamente através da integração entre os sistemas
físico e digital \cite{DIHAN2024e26503}.
Gêmeos Digitais aplicados à gestão urbana começaram a surgir por volta de 2018, potencializados
pelo conceito das \textit{smart cities}, onde cidadãos e a administração pública cooperam em prol
de tornar a cidade mais eficiente, inteligente, sustentável, segura, inclusiva e democrática.

Os Gêmeos Digitais Urbanos (GDU) são relativamente recentes; porém, vale destacar que não são a
primeira tecnologia aplicada à geografia urbana, que, na década de 2000, já contava com o
\textit{Building Information Modeling} (BIM) e o \textit{Geographic Information System} (GIS), os
quais são modelagens 3D das edificações e do espaço externo.
A diferença fundamental entre estes e o GDU reside no fluxo de dados: enquanto o BIM e o GIS são
modelos estáticos, o Gêmeo Digital é dinâmico e projetado para atuar com tecnologias de IoT,
refletindo assim sempre informações do mundo real, como temperatura, ruído sonoro, ocupação, etc.

Contudo, ainda há um problema latente de falta de padronização e de interoperabilidade na
implementação dos GDUs, visto que poucos deles conseguem compartilhar dados com outros GDUs ou
organizações \cite{FERREBIGORRA2022103905}.
Dessa demanda, em 2008, foi desenvolvido pelo \textit{Open Geospatial Consortium} (OGC) o
\textit{City Geography Markup Language} (CityGML), um padrão internacional para representar e
trocar dados de modelos 3D de cidades.
Este já é usado em regiões da Europa, sendo que todas as capitais europeias disponibilizam modelos
CityGML, bem como alguns países na Ásia e na África \cite{buildings13071758}.

O CityGML representa um grande passo na padronização da troca de informação entre sistemas de GDUs.
Todavia, segundo \citeonline[p.~103, tradução nossa]{metral_ontologybased_2010}, "[.
  ..] desenvolver um modelo urbano universal único é impossível; a percepção da realidade urbana é diversa e múltipla".
Por isso, propõe o uso de ontologias formais como ponte semântica entre diferentes modelos urbanos,
enquanto mantém o CityGML como modelo genérico para conceitos centrais compartilhados por
diferentes representações especializadas de dados urbanos.

As ontologias são definições formais de instâncias de objetos, grupos de objetos, e as relações
entre estes dentro de um domínio.
São representadas usando \textit{Web Ontology Language} (OWL) e \textit{Resource Description
  Framework} (RDF), e, diferente de esquemas de bancos de dados, descrevem o mundo real, não a
maneira como a informação é armazenada \cite{metral_ontologybased_2010}.
O uso de ontologias dentro de um mesmo contexto (e.g., cidade ou campus universitário) pode
enfrentar desafios semânticos quando temos fontes de dados heterogêneas, como BIM, GIS, sensores de
IoT, etc. Então, sob esse olhar, a abordagem de \citeonline{metral_ontologybased_2010} mostra-se
ainda mais conveniente, dada a possibilidade de integrar estruturas de dados diversas, definindo
relações de compatibilidade semântica entre uma ontologia específica e outra \cite{SHI2023102114}.

Por fim, a ontologia \textit{Sensor, Observation, Sample, and Actuator/Semantic Sensor Network}
(SOSA/SSN), proposta para sistemas de IoT, complementa o CityGML na camada de sensores.
Oriunda do trabalho em conjunto do OGC com o \textit{World Wide Web Consortium} (W3C), a primeira,
SOSA, define os conceitos de sensor, observação, amostragem e atuador, enquanto a SSN a estende,
incluindo também instruções de procedimentos, propriedades observadas e sistemas de sensores
\cite{haller2019}; Isso permite que dispositivos distintos, como uma câmera de segurança e um
detector de presença, observando um evento físico com dialetos diferentes, cheguem à mesma
conclusão (e.g. "há uma pessoa no pátio").
Assim, CityGML, SOSA/SSN e OWL/RDF formam os pilares semânticos para a modelagem do campus
Higienópolis do Instituto Presbiteriano Mackenzie, objeto central deste trabalho.

\section{Metodologia}

Este trabalho é, quanto ao tipo, uma pesquisa aplicada, visto que busca gerar insumos para resolver
uma dificuldade real: a falta de homogeneidade nas aplicações de GDU.
Quanto à abordagem, é formal, porque o objeto central de estudo é a modelagem formal de uma
ontologia aplicada ao contexto do campus da UPM.
Em relação aos objetivos, é de natureza exploratória, pois a aplicação de ontologias OWL como
camada semântica integradora em contextos de campus universitário, como instância de GDU, ainda é
incipiente na literatura.
Por fim, quanto ao procedimento técnico, adota-se o \textit{Design Science Research} (DSR).

O DSR é adotado devido ao foco em criar e avaliar um artefato para resolver um problema complexo
e/ou prático, ao passo que contribui com novos conhecimentos \cite{hevner_2004}.
Com isso, este trabalho deve gerar um artefato de modelo, que é a ontologia para representar
formalmente as entidades presentes no campus, de maneira alinhada aos padrões CityGML e SOSA/SSN.
Portanto, o conhecimento científico emerge do processo de construção, refinamento e validação do
próprio artefato frente à problemática envolvida.

A avaliação do artefato produzido terá como critério permitir uma maneira de navegar nas ontologias
estudadas, e, usando de motores como o \textit{SPARQL Protocol and RDF Query Language} (SPARQL),
responder consultas semânticas simples, como: "quais prédios estão paralelos à Rua da Consolação?
"; "quais espécies de árvores estão presentes no Bosque das Fadas?"; ou, ainda, "quais sensores observam a entrada do prédio 31?".

\section{Cronograma de Atividades}

\begin{table}[htbp]
  \centering
  \begin{tabular}{l|c|c|c|c|c}
    Atividade                          & Ago.
                                       & Set. & Out. & Nov. & Dez.     \\
    \hline
    Levantamento Físico do Campus      & X    &      &      &      &   \\
    Mapeamento de Entidades Principais & X    & X    &      &      &   \\
    \textit{Reasoner} e SPARQL         &      & X    &      &      &   \\
    Implementação da Arquitetura       &      & X    & X    & X    &   \\
    Resultados e Discussão             &      &      &      & X    & X \\
    Conclusão e Revisão                &      &      &      &      & X \\
  \end{tabular}
  \caption{Cronograma de atividades a serem desenvolvidas.}
  \label{tab:bullets}
\end{table}

% \section{Resultados e Discussão}

% Resultados e discussão ainda a serem esquematizados e apresentados mediante a implementação feita na fase de TCC II.

% \section{Conclusão}

% Conclusão mediante a implementação feita na fase de TCC II.

% \section{Referências bibliográficas}
% trabalhos não referenciados explicitamente no texto, mas consultados
% \nocite{SHI2023102114}

% ----------------------------------------------------------
% ELEMENTOS PÓS-TEXTUAIS
% ----------------------------------------------------------
\postextual

% ----------------------------------------------------------
% Referências bibliográficas
% ----------------------------------------------------------
\bibliography{referencias}

% ----------------------------------------------------------
\end{document}
